\documentclass{article}
\usepackage{amsmath,amssymb}
\usepackage{xurl}
\usepackage[colorlinks=true,linkcolor=blue,citecolor=blue,urlcolor=blue]{hyperref}
% Shared commands for notes in docs/paper-gaps/.

\DeclareUnicodeCharacter{2080}{\ifmmode{}_{0}\else\textsubscript{0}\fi}
\DeclareUnicodeCharacter{2081}{\ifmmode{}_{1}\else\textsubscript{1}\fi}
\DeclareUnicodeCharacter{2082}{\ifmmode{}_{2}\else\textsubscript{2}\fi}
\DeclareUnicodeCharacter{2099}{\ifmmode{}_{n}\else\textsubscript{n}\fi}

\newcommand{\Lean}{\textsc{Lean}}
\ExplSyntaxOn
\NewDocumentCommand{\leanid}{m}{
  \ifmmode
    \textsf{\detokenize{#1}}
  \else
    \group_begin:
    \urlstyle{sf}
    \tl_set:Nn \l_tmpa_tl {#1}
    \tl_replace_all:Nnn \l_tmpa_tl {\_} {_}
    \exp_args:NV \path \l_tmpa_tl
    \group_end:
  \fi
}
\ExplSyntaxOff
\providecommand{\tr}{\operatorname{tr}}
\newcommand{\tnleanIssueBase}{https://github.com/LionSR/TNLean/issues/}
\newcommand{\tnleanPrBase}{https://github.com/LionSR/TNLean/pull/}
\newcommand{\tnleanDocBase}{https://sirui-lu.com/TNLean/docs/}
\newcommand{\ghissue}[1]{\href{\tnleanIssueBase#1}{GitHub issue \##1}}
\newcommand{\ghpr}[1]{\href{\tnleanPrBase#1}{PR \##1}}
\newcommand{\doclink}[1]{\href{\tnleanDocBase#1.html}{\path{#1}}}

% Dirac notation and the identity operator, as in blueprint/src/macros/common.tex.
% \providecommand keeps note-local definitions authoritative.
\usepackage{dsfont}
\providecommand{\ket}[1]{|#1\rangle}
\providecommand{\bra}[1]{\langle #1|}
\providecommand{\braket}[2]{\langle #1 | #2 \rangle}
\providecommand{\ketbra}[2]{|#1\rangle\!\langle #2|}
\providecommand{\Id}{\mathds{1}}
\providecommand{\one}{\mathds{1}}
\providecommand{\C}{\mathbb{C}}
\providecommand{\MN}[1]{M_{#1}(\C)}
\providecommand{\MD}{M_D(\C)}

% Verdict marker: \gapnote{<kind>}{<status>}. Kind is the policy
% classification (clarification, local-correction, scope-restriction,
% unfaithful, false-source, open-gap); status is open, wip (elimination
% underway), resolved, or historical. Machine-read by the site index;
% prints a verdict line.
\newcommand{\gapnote}[2]{\par\noindent\textbf{Verdict:} #1 (#2).\par}

\providecommand{\leanid}[1]{\textsf{\detokenize{#1}}}
\providecommand{\gapnote}[2]{\par\noindent\textbf{Verdict:} #1 (#2).\par}
\title{Birth-Template Separation for Replacements Only}
\author{}
\date{October 15, 2026}
\begin{document}
\maketitle
\gapnote{scope-restriction}{open}

\section{The source statement}

Lemma~11.5 of~\cite{OpenAI2026AreaLaw} (\path{10-geometry.tex}, lines~573--582)
asserts a fixed constant $c_1>0$ with the following property. Let
$\mathcal B$ be the closed birth region either of an initial belt tile of side
scale $u=t_k$, or of a replacement tile whose removal square has half-side $u$.
Then $\mathcal B$ has distance at least $c_1u^{\rho_0}$ from the final actual
region of every earlier or peer identifier of the same label. Its proof,
lines~584--659, compares closed births with the older and peer regions after the
removal step, uses a compactness argument away from contact marks and an angular
estimate near them, and lets the following trim remove a square of radius
comparable with $u^{\rho_0}$ about every contact mark.

\section{The formal statement}

The formalization proves the replacement half of the lemma, for the actual
repair stages. Its only premise beyond the choice of the stopping scale is the
conclusion of the uniform repair induction, Lemma~11.4 (lines~471--485), stated
as a separate proposition.\footnote{\leanid{TNLean.PEPS.AreaLaw.Geometry.IsRepairBirthSeparated},
\leanid{TNLean.PEPS.AreaLaw.Geometry.exists_isRepairBirthSeparated_of_uniformRepairInductionStatement}.}
The constant $c_1$ is fixed before the repair data. The separation is proved
against every other identifier of the same generation with the same label,
which contains the earlier and peer identifiers of the source.

The births of initial belt tiles are not covered. For them the comparison is
with the primary and belt closures at the initial stage; it requires the local
classification of the initial belt arrangement around a reference cell rather
than the fixed replacement fan used here.

\section{What remains}

The initial-belt half of Lemma~11.5 remains to be formalized. The normalized
compactness and angular estimates already proved apply unchanged once the
initial belt closures near a reference cell are written as finite unions of
integer half-plane polygons, with the canonical sector assignment at the
reference cell marks supplying the contacts.

\bibliographystyle{alpha}
\bibliography{references}

\end{document}
